\begin{tikzpicture}[
    auto,
    font=\footnotesize\sffamily,
    node distance=1.1cm and 1.8cm,
    splitstate/.style 2 args={
      rectangle, rounded corners=4pt, minimum width=4.1cm, minimum height=1.35cm,
      align=center, draw=black!80, line width=0.8pt,
      path picture={
        \fill[#1] (path picture bounding box.north west) rectangle (path picture bounding box.east);
        \fill[#2] (path picture bounding box.west) rectangle (path picture bounding box.south east);
        \draw[black!20, dashed] (path picture bounding box.west) -- (path picture bounding box.east);
      }
    },
    intermitstate/.style 2 args={
      rectangle, rounded corners=4pt, minimum width=4.1cm, minimum height=1.35cm,
      align=center, draw=orange!85!black, line width=0.9pt,
      path picture={
        \fill[#1] (path picture bounding box.north west) rectangle (path picture bounding box.east);
        \fill[#2] (path picture bounding box.west) rectangle (path picture bounding box.south east);
        \draw[black!20, dashed] (path picture bounding box.west) -- (path picture bounding box.east);
      }
    },
    startnode/.style={circle, minimum size=0.85cm, align=center, draw=black!70, fill=blue!10, thick},
    arrow/.style={-{Stealth[length=2.5mm]}, line width=0.9pt, draw=black!75},
    cfgarrow/.style={-{Stealth[length=2.5mm]}, line width=0.85pt, draw=orange!80!black, dashed}
]

  % Nodo de Inicio y Arranque Seguro
  \node (start) [startnode] {Inicio};
  
  \node (boot) [splitstate={red!20}{red!20}, below=0.7cm of start] {
    \textbf{Arranque Seguro (2 s)}\\
    NS: Rojo \quad EO: Rojo\\
    \scriptsize Despeje inicial (Estado 2)
  };

  % Columna 1 (Izquierda): Fase 1 (Norte-Sur vehicular activa)
  \node (st0) [splitstate={green!22}{red!20}, below=0.9cm of boot] {
    \textbf{Estado 0}\\
    NS: Verde (+ Peatonal)\\
    EO: Rojo\\
    \scriptsize Paso vehicular Norte--Sur
  };

  \node (st1) [splitstate={yellow!35}{red!20}, below=1.0cm of st0] {
    \textbf{Estado 1}\\
    NS: Amarillo\\
    EO: Rojo\\
    \scriptsize Advertencia fin de paso NS
  };

  \node (st2) [splitstate={red!20}{red!20}, below=1.0cm of st1] {
    \textbf{Estado 2}\\
    NS: Rojo\\
    EO: Rojo\\
    \scriptsize Despeje de encrucijada
  };

  % Columna 2 (Centro): Fase 2 (Este-Oeste vehicular activa)
  \node (st3) [splitstate={red!20}{green!22}, right=1.8cm of st2] {
    \textbf{Estado 3}\\
    NS: Rojo\\
    EO: Verde (+ Peatonal)\\
    \scriptsize Paso vehicular Este--Oeste
  };

  \node (st4) [splitstate={red!20}{yellow!35}, above=1.0cm of st3] {
    \textbf{Estado 4}\\
    NS: Rojo\\
    EO: Amarillo\\
    \scriptsize Advertencia fin de paso EO
  };

  \node (st5) [splitstate={red!20}{red!20}, above=1.0cm of st4] {
    \textbf{Estado 5}\\
    NS: Rojo\\
    EO: Rojo\\
    \scriptsize Despeje de encrucijada
  };

  % Columna 3 (Derecha): Régimen Intermitente / Destellante (Estados 6 y 7)
  \node (st6) [intermitstate={yellow!35}{yellow!35}, right=1.8cm of st5] {
    \textbf{Estado 6}\\
    NS: Amarillo\\
    EO: Amarillo\\
    \scriptsize Intermitente ON (Precaución)
  };

  \node (st7) [intermitstate={black!10}{black!10}, below=1.2cm of st6] {
    \textbf{Estado 7}\\
    NS: Apagado\\
    EO: Apagado\\
    \scriptsize Intermitente OFF (Cadencia)
  };

  % Conexiones del flujo de arranque y ciclo vehicular
  \draw [arrow] (start) -- (boot);
  \draw [arrow] (boot) -- (st0);
  \draw [arrow] (st0) -- node [right, scale=0.75] {Vence Int 1} (st1);
  \draw [arrow] (st1) -- node [right, scale=0.75] {Vence Int 2} (st2);
  \draw [arrow] (st2) -- node [above, scale=0.75] {Vence Int 3} (st3);
  \draw [arrow] (st3) -- node [right, scale=0.75] {Vence Int 4} (st4);
  \draw [arrow] (st4) -- node [right, scale=0.75] {Vence Int 5} (st5);

  % Retorno directo del ciclo vehicular cuando Int 7 = 0 y Int 8 = 0
  \draw [arrow] (st5.west) -- node [above, scale=0.72, align=center] {Vence Int 6\\{\scriptsize (si Int 7, 8 = 0 s)}} (st0.east);

  % Conexión hacia Estados 6 y 7 cuando Int 7 > 0 s
  \draw [cfgarrow] (st5.east) -- node [above, scale=0.72, align=center] {Int 7 $>$ 0 s} (st6.west);
  \draw [arrow] (st6) -- node [right, scale=0.75] {Vence Int 7} (st7);

  % Retorno desde Estado 7 al inicio del ciclo (perimetral inferior) o lazo intermitente
  \coordinate (under) at ($(st2.south)+(0,-0.55cm)$);
  \draw [cfgarrow] (st7.south) -- (st7.south |- under) 
    -- node [below, scale=0.72] {Vence Int 8 $\to$ reinicia ciclo en Estado 0}
    ($(st0.west |- under)+(-0.55cm,0)$) 
    -- ($(st0.west)+(-0.55cm,0)$) -- (st0.west);
  \draw [cfgarrow] (st7.east) to[out=0, in=0] node [right, scale=0.72, align=left] {Si Int 1--6 = 0 s\\(Destello puro)} (st6.east);

\end{tikzpicture}
